\documentclass[11pt]{article}
\usepackage{amsmath, amssymb, amsfonts}
\usepackage{geometry}
\geometry{margin=2.5cm}

\title{Bedeutung der Abbildungen \texttt{perfektoid\_to\_lp} und \texttt{lp\_to\_perfektoid}}
\author{}
\date{}

\begin{document}
\maketitle

\section*{Kurzüberblick}

Es existieren zwei zentrale Abbildungen:
\[
\texttt{perfektoid\_to\_lp}
\quad\text{und}\quad
\texttt{lp\_to\_perfektoid}.
\]

\begin{itemize}
\item \textbf{\texttt{perfektoid\_to\_lp}} übersetzt eine endliche perfektoide (p-adische) Darstellung
      in ein lineares Optimierungsproblem, das vom Simplex-Algorithmus verarbeitet werden kann.
\item \textbf{\texttt{lp\_to\_perfektoid}} interpretiert eine Lösung dieses Optimierungsproblems
      wieder als perfektoide Darstellung.
\end{itemize}

Diese Abbildungen sind \emph{keine Inversen}, sondern bilden ein
\emph{Modellierungs-- und Interpretationspaar}.

\section*{1. Mathematische Ausgangsstruktur}

Eine endliche p-adische (perfektoide) Approximation wird beschrieben durch
\[
x \approx \sum_{i=0}^{n} a_i \, p^{k_0 + i},
\]
wobei
\[
p \text{ eine Primzahl}, \quad
k_0 \in \mathbb{Z}, \quad
a_i \in \{0,\dots,p-1\}.
\]

Diese Daten werden zusammengefasst als
\[
\text{PerfektoidApprox} = (p, k_0, (a_0,\dots,a_n)).
\]

\section*{2. Bedeutung von \texttt{perfektoid\_to\_lp}}

Die Abbildung
\[
\texttt{perfektoid\_to\_lp} :
(p, k_0, a_0,\dots,a_n)
\longmapsto
\text{LinearProgram}
\]
übersetzt arithmetische Struktur in Optimierungsstruktur.

\subsection*{2.1 Variablen}

Für jede p-adische Ziffer wird eine LP-Variable eingeführt:
\[
a_i \longleftrightarrow x_i.
\]

\subsection*{2.2 Zielfunktion}

Die Zielfunktion minimiert die p-adische Bewertung:
\[
\min \sum_{i=0}^{n} (k_0 + i)\, x_i.
\]

Dies bevorzugt Lösungen mit niedriger Bewertung und modelliert p-adische
``Kleinheit''.

\subsection*{2.3 Nebenbedingungen}

\paragraph{Ziffernbeschränkung}
\[
0 \le x_i \le p-1.
\]

\paragraph{Frobenius-Kohärenz (perfektoide Struktur)}
\[
x_{i+p} - x_i = 0,
\quad \text{für alle zulässigen } i.
\]

Diese Bedingung modelliert diskret die Surjektivität des Frobenius
\[
x \mapsto x^p.
\]

\subsection*{2.4 Ergebnis}

Das Resultat ist ein lineares Programm
\[
\min c^\top x \quad \text{unter} \quad A x = b, \; l \le x \le u,
\]
welches direkt vom Simplex-Algorithmus behandelt werden kann.

\section*{3. Bedeutung von \texttt{lp\_to\_perfektoid}}

Die Abbildung
\[
\texttt{lp\_to\_perfektoid} :
x^\ast \longmapsto (p, k_0, a_0,\dots,a_n)
\]
interpretiert eine numerische Lösung wieder arithmetisch.

\subsection*{3.1 Eingabe}

Eine Simplex-Lösung liefert einen reellen Vektor
\[
x^\ast = (x_0^\ast, \dots, x_n^\ast).
\]

\subsection*{3.2 Diskretisierung}

Die kontinuierlichen Werte werden gerundet:
\[
a_i := \operatorname{round}(x_i^\ast).
\]

Dadurch entsteht wieder eine zulässige p-adische Ziffernfolge.

\subsection*{3.3 Ergebnis}

Die rekonstruierte perfektoide Approximation lautet
\[
x_{\text{neu}} = \sum_{i=0}^{n} a_i \, p^{k_0 + i}.
\]

\section*{4. Warum dies keine echte Inversion ist}

Im Allgemeinen gilt:
\[
\texttt{lp\_to\_perfektoid} \circ \texttt{perfektoid\_to\_lp}
\neq \mathrm{id}.
\]

Gründe:
\begin{itemize}
\item Relaxierung von Ganzzahligkeit (LP statt MILP),
\item Rundungsfehler,
\item endliche Trunkierung,
\item linearisierte Frobenius-Bedingungen.
\end{itemize}

Die Abbildungen bilden daher keinen Isomorphismus, sondern einen
\emph{Modellierungszyklus}.

\section*{5. Intuitive Zusammenfassung}

\[
\begin{array}{ccl}
\text{p-adische Zahl} & \longrightarrow & \text{Optimierungsproblem} \\
\text{Bewertung}     & \longrightarrow & \text{Zielfunktion} \\
\text{Frobenius}     & \longrightarrow & \text{Nebenbedingungen} \\
\text{Lösung}        & \longrightarrow & \text{neue Ziffernfolge}
\end{array}
\]

\section*{6. Fazit}

\texttt{perfektoid\_to\_lp} kodiert perfektoide Struktur als lineare Geometrie,  
\texttt{lp\_to\_perfektoid} dekodiert numerische Lösungen wieder als Zahlen.

Optimiert wird nicht die Zahl selbst, sondern ihre
arithmetisch motivierte Darstellung.

\end{document}

